\documentclass[11pt,a4paper]{article}

\input{../../sablony/pripravy/priprava.sty}

\nastavPripravu{4. ročník}{Kombinatorika}{42}{Smíšené kombinatorické úlohy}

\begin{document}

\blok{Motivace a návaznost}

Ke každému zadání musí žák nejprve říct, co tvoří jeden výsledek: pořadí všech prvků, obsazení rolí, výběr skupiny, opakování s pevnými četnostmi, nebo opakované volby do kódu.

\poznamka{Předpoklady: celý blok 29--41. Cíl: samostatně zvolit model, řešit omezení a vysvětlit úplnost zvoleného postupu.}

\blok{Výklad a odvození}

\begin{definiceBox}
\textbf{Rozhodovací postup:} (1) Co přesně tvoří výsledek? (2) Záleží na pořadí? (3) Používáme všechny prvky? (4) Smí se opakovat a mají druhy pevné četnosti? (5) Jaká omezení musí splňovat začátek, konec a sousednost?
\end{definiceBox}
Nejdřív vysvětli zvolenou metodu. Následně spočítej výsledek a ověř, že nebyl započten vícekrát. U velkých faktoriálů může zůstat odpověď v přesném faktoriálovém tvaru.

\blok{Společně řešené příklady}

\begin{prikladReseniBox}[title={Příklad 1}]
Šest různých lidí se řadí do fronty. Dvě předem určené osoby musí stát vedle sebe. Kolik pořadí je možných?
\tcblower
Spojíme dvojici do bloku: $2\cdot5!=240$.
\end{prikladReseniBox}



\begin{prikladReseniBox}[title={Příklad 2}]
Z osmi lidí vybíráme tříčlenný tým bez funkcí. Kolik možností? Kolik, jestliže obsazujeme tři různé funkce?
\tcblower
Tým $\binom83=56$; funkce $V(8,3)=8\cdot7\cdot6=336$.
\end{prikladReseniBox}



\begin{prikladReseniBox}[title={Příklad 3}]
Kolik různých řetězců délky šest ze znaků A, A, B, B, B, C lze vytvořit?
\tcblower
$6!/(2!3!)=60$.
\end{prikladReseniBox}



\begin{prikladReseniBox}[title={Příklad 4}]
Kolik čtyřmístných kódů nad písmeny A, B, C, D existuje s opakováním?
\tcblower
$4^4=256$.
\end{prikladReseniBox}



\blok{Samostatná práce}

\begin{zadaniReseniBox}[title={Samostatná práce 1}]
Z desetimístné kandidátky vybíráme čtyřčlenný výbor bez funkcí.
\tcblower
$\binom{10}{4}=210$.
\end{zadaniReseniBox}



\begin{zadaniReseniBox}[title={Samostatná práce 2}]
Kolik různých přeskupení slova KOKOS vznikne?
\tcblower
Dvě K, dvě O, jedno S: $5!/(2!2!)=30$.
\end{zadaniReseniBox}



\begin{zadaniReseniBox}[title={Samostatná práce 3}]
Kolik pětimístných PINů z číslic 0 až 9 existuje, nejsou-li dovoleny stejné sousední cifry?
\tcblower
$10\cdot9^4=65610$.
\end{zadaniReseniBox}



\blok{Rozšíření}

\begin{zadaniReseniBox}[title={Rozšíření 1}]
Na fotografii stojí osm různých postav. A musí být před B a C musí stát poslední. Kolik pořadí vyhovuje?
\tcblower
Poslední místo obsadí C; ze zbylých sedmi postav je A před B právě v polovině pořadí: $7!/2=2520$.
\end{zadaniReseniBox}



\begin{zadaniReseniBox}[title={Rozšíření 2}]
\textbf{Volitelné:} Kolika způsoby lze z pěti příchutí vybrat čtyři kopečky, když se chuť může opakovat a pořadí nerozlišujeme?
\tcblower
Kombinace s opakováním: $\binom{5+4-1}{4}=\binom84=70$. Pouze pro třídu, kde byl tento typ vysvětlen.
\end{zadaniReseniBox}



\blok{Domácí úkol}

\poznamka{Doporučená práce doma: přibližně 15--30 minut. Úlohy jsou seřazeny od základních k náročnějším; učitel může poslední podúlohu označit jako volitelnou.}

\begin{zadaniReseniBox}[title={Domácí úkol}]

\begin{enumerate}[label=\alph*),leftmargin=*,itemsep=2pt]

\item Ze sedmi studentů se seřadí všichni; Petra musí být první. Urči počet.
\item Ze sedmi studentů vybereme předsedu a zapisovatele. Urči počet.
\item Ze sedmi studentů vybereme dvojčlennou komisi bez rolí. Urči počet.
\item Kolik různých přeskupení má slovo ABABA?
\item Kolik čtyřmístných PINů z číslic $0,1,2,3$ existuje s opakováním?
\item Kolik tříčlenných komisí ze šesti lidí obsahuje Annu?
\item Osm různých lidí se seřadí do fronty. Dvě konkrétní osoby musejí stát vedle sebe. Urči počet pořadí.

\end{enumerate}

\tcblower

\begin{enumerate}[label=\alph*),leftmargin=*,itemsep=2pt]

\item $6!=720$ (permutace).
\item $7\cdot6=42$ (variace).
\item $\binom72=21$ (kombinace).
\item $5!/(3!2!)=10$ (permutace s opakováním).
\item $4^4=256$ (variace s opakováním).
\item Anna je pevně, vybereme dvě z pěti: $\binom52=10$.
\item $2\cdot7!=10080$.

\end{enumerate}

\end{zadaniReseniBox}

\poznamka{Orientační práce v hodině: 0--10 min klasifikační tabulka; 10--32 min smíšené úlohy; 32--42 min samostatná práce; 42--45 min shrnutí.}

\poznamka{Domácí práce je současně individuální přípravou na krátký průběžný test. Nehodnoť pouhý výběr vzorce bez slovního zdůvodnění.}

\end{document}
